% !TeX root = ../复分析笔记.tex
% 由主文档通过 \input 载入；请编译主文档。
\setcounter{chapter}{11} % 保留原书章号。
\chapter{调和函数与圆盘边值问题}
\label{chap:ca-harmonic}

\section{从解析函数的实部到拉普拉斯方程}

\begin{definition}[调和函数与调和共轭]
实值函数 $u\in C^2(D)$ 若满足
\begin{equation}
 \Delta u=u_{xx}+u_{yy}=0,
 \label{eq:ca-laplace}
\end{equation}
称为区域 $D$ 上的\emph{调和函数（harmonic function）}。若 $u+iv$ 全纯，则称 $v$ 是 $u$ 的\emph{调和共轭（harmonic conjugate）}。
\end{definition}

在理想的均匀介质、无内部热源的二维稳态导热模型中，温度可满足拉普拉斯方程；在无电荷的平面区域内，静电势也可满足同样的方程。这些是带物理假设的模型解释。数学上，本章研究的是给定方程、区域及边界数据后，解是否存在、是否唯一、怎样表示。

由 C--R 方程和全纯函数的光滑性，全纯函数的实部与虚部都调和。但反方向需更细致地表述：一个调和函数局部可以成为全纯函数的实部，能否在整个区域上找到单值调和共轭，取决于区域的拓扑。

\begin{theorem}[调和共轭的存在]
\label{thm:ca-conjugate}
每个调和函数在每一点的小邻域内都有调和共轭。若定义域 $D$ 单连通，则存在全局调和共轭；在连通区域上，共轭函数在相差实常数的意义下唯一。
\end{theorem}
\begin{proof}
由 C--R 方程，应要求 $v_x=-u_y$、$v_y=u_x$。在以 $(x_0,y_0)$ 为基点的小矩形内定义
\[
 v(x,y)=-\int_{x_0}^{x}u_y(s,y_0)\,\mathrm ds
 +\int_{y_0}^{y}u_x(x,t)\,\mathrm dt.
\]
直接微分得 $v_y=u_x$，又由 $u_{xx}=-u_{yy}$，有
\[
 v_x=-u_y(x,y_0)+\int_{y_0}^{y}u_{xx}(x,t)\,\mathrm dt
 =-u_y(x,y).
\]
因此 $u+iv$ 全纯，得到局部存在性。对单连通区域，可沿路径定义
\begin{equation}
 v(z)=v(z_0)+\int_{z_0}^{z}(-u_y\,\mathrm dx+u_x\,\mathrm dy).
 \label{eq:ca-conjugate-integral}
\end{equation}
这个实微分形式的闭路积分为零：局部路径无关性由刚才的构造给出，再沿闭曲线收缩过程拼接局部原函数，方法与定理~\ref{thm:ca-cauchy-domain} 相同；对围住普通平面区域的简单闭曲线，也可直接用格林公式把积分化成 $\iint\Delta u\,\mathrm dx\mathrm dy=0$。故全局定义良好。若 $v_1,v_2$ 都满足要求，则其差的两个偏导均为零，连通性使差为常数。
\end{proof}

\begin{example}[实际求调和共轭]
设 $u=x^2-y^2+2x$，求 $v$ 使 $f=u+iv$ 在整个平面全纯。
\end{example}
\begin{solve}
先检查 $u_{xx}+u_{yy}=2-2=0$。由 $v_y=u_x=2x+2$，积分得 $v=(2x+2)y+\phi(x)$。再由 $v_x=-u_y=2y$，得到 $\phi'(x)=0$。所以 $v=2xy+2y+C$，相应 $f(z)=z^2+2z+iC$，$C\in\mathbb R$。不先检查调和性，也可以在第二个方程中发现不相容，但先检查可避免无效积分。
\end{solve}

在去心平面上，$u(z)=\ln|z|$ 调和，却没有全局单值调和共轭。其候选共轭是辐角，沿单位圆计算式~\eqref{eq:ca-conjugate-integral} 的增量为 $2\pi$，不满足路径无关性。删去一条支割线后，则可取 $v=\operatorname{Arg}z$。这把对数分支问题与调和共轭问题联系起来。

\section{平均值性质与极值原理}

\begin{theorem}[平均值性质]
\label{thm:ca-harmonic-mean}
若 $u$ 在包含闭圆盘 $\overline{B(a,r)}$ 的开集上调和，则
\begin{equation}
 u(a)=\frac1{2\pi}\int_0^{2\pi}u(a+re^{it})\,\mathrm dt
 =\frac1{\pi r^2}\iint_{B(a,r)}u\,\mathrm dx\mathrm dy.
 \label{eq:ca-harmonic-mean}
\end{equation}
\end{theorem}
\begin{proof}
在比该闭圆盘略大的圆盘内取调和共轭，得到全纯函数 $F$ 且 $u=\operatorname{Re}F$。柯西积分公式经圆周参数化给出 $F(a)$ 为圆周平均值，取实部得第一式。对每个 $0<s<r$ 使用圆周平均值，再乘 $2\pi s$ 对 $s$ 积分，利用极坐标面积公式得第二式。
\end{proof}

平均值性质表示中心值既不偏高也不偏低，而是周围各方向取值的平衡。它导致调和函数的极值不能孤立地出现在内部。

\begin{theorem}[调和函数的最大、最小值原理]
\label{thm:ca-harmonic-max}
若调和函数 $u$ 在区域 $D$ 内取得其在 $D$ 上的最大值或最小值，则 $u$ 为常数。若 $D$ 有界且 $u\in C(\overline D)$、在 $D$ 调和，则最大值和最小值都可在边界取得。因此同一连续边界数据至多对应一个这样的调和函数。
\end{theorem}
\begin{proof}
设最大值为 $M$，在 $a\in D$ 取得。对任意足够小的圆，其上 $u\le M$，平均值又等于 $M$，由连续性可知圆上每一点都等于 $M$。改变半径得 $u$ 在 $a$ 附近恒为 $M$。于是集合 $\{u=M\}$ 在 $D$ 中非空、开且闭，由连通性等于 $D$。对 $-u$ 应用同样论证得最小值结论。有界区域闭包紧致，极值存在；若两个解边界相同，其差边界为零，用最大、最小值原理得差处处为零。
\end{proof}

\section{泊松公式与圆盘上的狄利克雷问题}

所谓\emph{狄利克雷问题（Dirichlet problem）}，是给定边界值，求区域内部满足拉普拉斯方程并连续达到该边界值的函数。先在单位圆盘上讨论，设边界函数 $\varphi(t)$ 为连续、实值、$2\pi$ 周期函数。

\begin{theorem}[圆盘上的泊松公式]
\label{thm:ca-poisson}
对 $0\le r<1$，定义
\begin{equation}
 u(re^{i\theta})=\frac1{2\pi}\int_0^{2\pi}
 P_r(\theta-t)\varphi(t)\,\mathrm dt,
 \qquad
 P_r(s)=\frac{1-r^2}{1-2r\cos s+r^2}.
 \label{eq:ca-poisson}
\end{equation}
则 $u$ 在单位圆盘调和，并可连续延拓到闭圆盘，边界值为 $u(e^{it})=\varphi(t)$。它是该边值问题的唯一解。$P_r$ 称为\emph{泊松核（Poisson kernel）}。
\end{theorem}
\begin{proof}
首先注意
\[
 P_r(\theta-t)=\operatorname{Re}\frac{e^{it}+z}{e^{it}-z},\qquad z=re^{i\theta}.
\]
因此 $u$ 是全纯函数
\[
 F(z)=\frac1{2\pi}\int_0^{2\pi}\frac{e^{it}+z}{e^{it}-z}\varphi(t)\,\mathrm dt
\]
的实部。对每个较小闭圆盘，分母统一远离零，核及其导数有统一界，可对积分取差商并交换极限，所以 $F$ 全纯，$u$ 调和。

核具有三个性质：$P_r(s)>0$；其一个周期上的积分为 $2\pi$；对任意固定角距离 $\delta>0$，在 $\delta\le|s|\le\pi$ 上 $P_r(s)$ 随 $r\uparrow1$ 一致趋于零。第二条可由几何级数得到
\[
 P_r(s)=1+2\sum_{n=1}^\infty r^n\cos(ns),\qquad r<1,
\]
再逐项积分；第三条来自分母远离零。于是
\[
 u(re^{i\theta})-\varphi(\theta)
 =\frac1{2\pi}\int_0^{2\pi}P_r(\theta-t)
 [\varphi(t)-\varphi(\theta)]\,\mathrm dt.
\]
靠近 $t=\theta$ 的部分由 $\varphi$ 的一致连续性控制，远离部分由第三条控制，得到误差对 $\theta$ 一致趋于零。结合 $\varphi$ 连续，便得到沿任意方向趋近边界时的连续延拓。唯一性由调和极值原理给出。
\end{proof}

泊松核把边界值做加权平均；$r=0$ 时所有方向权重相同，越靠近边界，权重越集中在附近的边界点。若半径为 $R$、圆心为 $a$，通过缩放 $z=a+R\zeta$ 可得
\[
 u(a+re^{i\theta})=\frac1{2\pi}\int_0^{2\pi}
 \frac{R^2-r^2}{R^2-2Rr\cos(\theta-t)+r^2}
 \varphi(t)\,\mathrm dt,\qquad 0\le r<R,
\]
其中 $\varphi(t)$ 是点 $a+Re^{it}$ 上的边界值。

\begin{example}[由边界的三角多项式求内部解]
求单位圆盘内的调和函数，其边界值为 $\varphi(t)=3+2\cos t-\sin(2t)$。
\end{example}
\begin{solve}
由泊松核展开及三角函数正交性，常数项不变，第 $n$ 个频率乘以 $r^n$，所以
\[
 u(re^{i\theta})=3+2r\cos\theta-r^2\sin(2\theta)
 =3+2x-2xy.
\]
直接求拉普拉斯算子得到零，代入单位圆得到给定边界值，再由唯一性确定这是所求解。一般而言，若边界是有限三角多项式，逐频率乘以 $r^n$ 是最便捷的计算方式；对任意连续边界数据，泊松积分公式始终适用，不需要先假设其傅里叶级数逐点收敛。
\end{solve}

\begin{proposition}[平均值性质刻画调和性]
\label{prop:ca-mean-characterization}
若实值连续函数 $u$ 在区域 $D$ 上满足：对每个 $\overline{B(a,r)}\subset D$ 都有 $u(a)=\frac1{2\pi}\int_0^{2\pi}u(a+re^{it})\,\mathrm dt$，则 $u$ 在 $D$ 调和。
\end{proposition}
\begin{proof}
任取闭圆盘 $\overline{B(a,R)}\subset D$。以 $u$ 的圆周值为边界数据作泊松积分，得到在闭圆盘连续、内部调和的 $v$。差 $h=u-v$ 连续，边界为零，且满足所有内部小圆的平均值性质。最大值原理的平均值证明只使用连续性与此性质：若正最大值在内部取得，所有充分小圆上的值必须相同，最大值集合因而开且闭；连通性迫使整个圆盘内恒等于该正数，与零边界矛盾。同理不可能有负最小值，故 $h=0$。于是 $u=v$ 在这个圆盘调和。每点均可包含于这样的圆盘，故 $u$ 在 $D$ 调和，特别地自动具有所需的二阶连续偏导。
\end{proof}

\section{正调和函数的估计与共形变换}

若 $u$ 在闭圆盘 $\overline{B(0,R)}$ 的邻域调和且非负，由泊松核的上下界，可得对 $|z|=r<R$，
\begin{equation}
 \frac{R-r}{R+r}u(0)\le u(z)\le\frac{R+r}{R-r}u(0).
 \label{eq:ca-harnack}
\end{equation}
这是圆盘上的哈纳克不等式（Harnack inequality）。证明只需观察泊松核的最小、最大值分别为 $(R-r)/(R+r)$ 与 $(R+r)/(R-r)$，再用边界平均值等于 $u(0)$。它表明正调和函数在远离边界的小范围内不能任意剧烈变化；靠近边界时常数变大，也提醒我们内部控制与边界行为有不同尺度。

\begin{proposition}[调和性的共形不变性]
若 $u$ 在区域 $G$ 调和，$f:D\to G$ 全纯，则 $u\circ f$ 在 $D$ 调和。若 $f$ 还是双全纯映射，则可以通过它在两个区域之间转换调和函数。
\end{proposition}
\begin{proof}
写 $f=p+iq$，对二阶偏导使用实链式法则。由 $p,q$ 调和以及 C--R 方程，交叉项抵消，得到
\[
 \Delta(u\circ f)=|f'(z)|^2(\Delta u)(f(z))=0.
\]
或者在每个小邻域写 $u=\operatorname{Re}F$，则 $u\circ f=\operatorname{Re}(F\circ f)$，也得到同一结论。
\end{proof}

\begin{example}[扇形上的边界值]
在扇形 $0<\arg z<\alpha$ 内，求两条边分别取 $0$、$1$ 的一个有界调和函数，其中边界条件只针对两条射线上的非零有限点。
\end{example}
\begin{solve}
选取 $0<\arg z<\alpha$ 的对数分支，取
\[
 u(z)=\frac{\arg z}{\alpha}=\operatorname{Im}\frac{\log z}{\alpha}.
\]
它在扇形内调和，值在 $(0,1)$，并达到指定射线边界值。原点处两侧数据冲突，不能要求那里连续；区域又无界，因此这里仅给出一个满足所述条件的解，不能未经额外论证直接套用有界区域的边值唯一性定理。
\end{solve}

\section{习题}
\label{sec:ca-harmonic-ex}
\begin{enumerate}
 \item 判断 $u=x^2+y^2$ 与 $v=e^x\cos y$ 是否调和；对调和者求一个调和共轭。
 \item 求 $u=x^3-3xy^2+2y$ 的调和共轭，并写出对应的全纯函数。
 \item 求单位圆盘内边界值为 $2+\cos(3t)+2\sin t$ 的调和函数，并写出中心值。
 \item 设 $u$ 在闭单位圆盘的邻域调和且边界满足 $|u|\le M$，证明圆盘内也有 $|u|\le M$。
 \item 证明 $\ln|z|$ 在原点外调和，并计算其候选共轭沿正向单位圆的增量。
\end{enumerate}

